\section{The local endomorphism algebra}
\label{sec:endomorphism-algebra}

The preceding results solve LU equivalence for every prime power.  This
section extracts additional structure from the same systematic matrix.  The
algebraic description below is valid over the underlying prime field; the
five-type classification and its code interpretations specialize explicitly
to prime local dimension.  None of this section is needed for the robust
rounding theorem.

Let \(m\geq2\), and let \(L\subseteq\bigoplus_iP_i\) be the label space of
a stabilizer \(\AME(2m,q)\) state, where each \(P_i\) is regarded as a
\(2e\)-dimensional \(\F_p\)-space.  Define
\[
 \mathcal E_L=
 \left\{(T_i)_i\in\prod_i\operatorname{End}_{\F_p}(P_i):
   \left(\bigoplus_iT_i\right)L\subseteq L\right\}.       \tag{5.1}
\]
Thus \(\mathcal E_L\) is the algebra of block-diagonal, one-party label
endomorphisms preserving \(L\).

\begin{theorem}[Endomorphism algebra from fundamental cycles]
\label{thm:local-endomorphism-algebra}
Fix a half-set \(B\sqcup B^c\), write \(G=K_B\), and choose
\(b\in B\) and \(c\in B^c\).  For \(i\ne b\) and \(j\ne c\), put
\[
 Q_{ij}=G_{bj}G_{ij}^{-1}G_{ic}G_{bc}^{-1}
   \in\operatorname{GL}_{2e}(\F_p).
\]
Evaluation at party \(b\) is an algebra isomorphism
\[
 \mathcal E_L\ \cong\
 \operatorname{Cent}_{\operatorname{End}_{\F_p}(P_b)}
       \{Q_{ij}:i\ne b,\ j\ne c\}.                    \tag{5.2}
\]
Under this isomorphism, the compatible local symplectic automorphisms are
the units whose propagated components preserve the symplectic form at every
party.  If \(q=p\), this condition is simply
\[
 \operatorname{Aut}_{\mathrm{tr}}(\psi)
 \cong\{T\in\mathcal E_L:\det T_b=1\}.               \tag{5.3}
\]
In particular, the abstract algebra does not depend on the chosen half-set
or base parties.
\end{theorem}

\begin{proof}
A block-diagonal endomorphism, written as \((A_i)_{i\in B}\) and
\((D_j)_{j\in B^c}\), preserves the graph of \(G\) exactly when
\[
 A_iG_{ij}=G_{ij}D_j\qquad(i\in B,\ j\in B^c).         \tag{5.4}
\]
Every block \(G_{ij}\) is invertible by
Proposition~\ref{prop:half-set-direct-sum}.  Once \(A_b=A\) is fixed, the
base row and column force
\[
 D_j=G_{bj}^{-1}AG_{bj},\qquad
 A_i=G_{ic}D_cG_{ic}^{-1}.                             \tag{5.5}
\]
The remaining equations are precisely \(AQ_{ij}=Q_{ij}A\).  Evaluation
at \(b\) is therefore bijective onto the common centralizer, and
\((5.5)\) shows that it respects addition and multiplication.

An automorphism is compatible exactly when every propagated component in
\((5.5)\) is symplectic.  For \(q=p\), all components are conjugate
\(2\times2\) matrices and \(\operatorname{Sp}_2(q)=\operatorname{SL}_2(q)\),
so this is equivalent to \(\det T_b=1\).  For \(e>1\), one must retain all
transported symplectic-form conditions; a single determinant condition is
not sufficient.  Finally, the definition in \((5.1)\) is independent of
the systematic presentation.
\end{proof}

For prime dimension, Theorem~\ref{thm:local-endomorphism-algebra} replaces
the former list of unrelated group orders by a classification of one
algebra.  It also explains what additional code structure each nonscalar
case supplies.

\begin{corollary}[Five prime-field algebra types]
\label{cor:endomorphism-types}
Assume \(m\geq2\) and \(q=p\), and put \(d=\gcd(2,q-1)\).  Exactly one row of
Table~\ref{tab:endomorphism-types} applies.  The middle two columns give
\(\operatorname{Aut}_{\mathrm{tr}}(\psi)\) and its order, and
\[
 |\Gamma_\psi|=q^{2m}|\operatorname{Aut}_{\mathrm{tr}}(\psi)|.
\]
After compatible local coordinates are chosen, the last column gives the
code or module structure forced by \(\mathcal E_L\).
\end{corollary}

\begin{table}[H]
\centering
\small
\renewcommand{\arraystretch}{1.15}
\begin{tabular}{@{}L{23mm}L{27mm}L{16mm}L{67mm}@{}}
\toprule
\(\mathcal E_L\) & compatible group & order & resulting structure \\
\midrule
\(M_2(\F_q)\) & \(\operatorname{SL}_2(q)\) & \(q(q^2-1)\) &
  \(\F_q^2\otimes C\), where \(C\) is a weighted self-dual classical MDS
  code; an MDS--CSS stabilizer after local symplectic coordinate changes \\
\(\F_q I\) & center of \(\operatorname{SL}_2(q)\) & \(d\) &
  no additional linear structure is forced \\
\(\F_q\times\F_q\) & split norm-one torus & \(q-1\) &
  two weighted-dual classical MDS codes on complementary local axes;
  an MDS--CSS stabilizer after local symplectic coordinate changes \\
\(\F_{q^2}\) & field norm-one torus & \(q+1\) &
  an \(\F_{q^2}\)-linear MDS code, weighted Hermitian self-dual in
  compatible coordinates \\
\(\F_q[\epsilon]/(\epsilon^2)\) & determinant-one units & \(qd\) &
  a free rank-\(m\) dual-number module for which every half-set projection
  is an isomorphism \\
\bottomrule
\end{tabular}
\caption{The five prime-field endomorphism algebras, their determinant-one
unit groups, and the code structures they force.}
\label{tab:endomorphism-types}
\end{table}

\begin{proof}[Proof of Corollary~\ref{cor:endomorphism-types}]
If every fundamental holonomy is scalar, its common centralizer is
\(M_2(\F_q)\).  Otherwise choose a nonscalar holonomy \(Q\).  Its
centralizer is the quadratic algebra \(\F_q[Q]\).  If another holonomy
does not commute with \(Q\), the common centralizer is scalar; if all
holonomies commute, the common centralizer remains \(\F_q[Q]\).  The three
quadratic possibilities are the split algebra, the quadratic field, and
the dual numbers.  The determinant restricts respectively to the product,
the field norm, and \(\det(aI+bN)=a^2\) for \(N^2=0\).  This gives the
five groups and orders in the table, including characteristic two.

The code interpretations follow from the action of \(\mathcal E_L\) on
the local planes and the fact that every half-set projection of \(L\) is
an isomorphism.  Appendix~\ref{app:endomorphism-code-structures} gives the
module argument.
\end{proof}

The classification is constructive without enumerating a finite group.
Indeed, \(\mathcal E_L\) is the kernel of the homogeneous linear system
\([A,Q_{ij}]=0\) in the four entries of \(A\).  Its dimension is
\(1\), \(2\), or \(4\); in the middle case the characteristic polynomial
of one nonscalar element distinguishes the three quadratic algebras.
Propagation by \((5.5)\) then recovers the one-party actions.  Thus the
recognition calculation also returns the symmetry type and the corresponding
code-structure row of Table~\ref{tab:endomorphism-types}.

The constituent CSS and Hermitian code constructions are standard; see,
for example, \cite{KetkarKlappeneckerKumarSarvepalli2006}.  In the qubit
case, invariance under a uniform order-three local Clifford was already
identified with \(\F_4\)-linearity
\cite[Theorem~2]{VanDenNestDehaeneDeMoor2005}.  Our use of the table is the
converse organization specific to stabilizer AME states: the intrinsic
endomorphism algebra determines both the compatible group and the code
structure.  We do not call the dual-number row an MDS code over a ring.

The algebra formulation also reveals a restriction at small party number
that is absent from the five-row list.

\begin{theorem}[Forced endomorphisms for four and six parties]
\label{thm:low-party-endomorphisms}
Let \(q=p\).  For every stabilizer \(\AME(2m,q)\) state with
\(m=2\) or \(m=3\),
\[
 \dim_{\F_q}\mathcal E_L\geq2.                         \tag{5.6}
\]
Equivalently, all fundamental four-cycle holonomies commute, so the scalar
row of Table~\ref{tab:endomorphism-types} cannot occur.  If \(q\geq5\),
the state consequently has a noncentral compatible local symplectic
automorphism.
\end{theorem}

\begin{proof}
For \(m=2\) there is only one fundamental holonomy, whose centralizer has
dimension at least two.

Let \(m=3\).  Independent invertible changes of block row and column bases
put the \(3\times3\) block matrix \(G\) in the form
\[
 N=\begin{pmatrix}
 I&I&I\\ I&A&B\\ I&C&D
 \end{pmatrix}.                                      \tag{5.7}
\]
These basis changes need not be symplectic.  In dimension two, however,
every invertible change rescales the alternating form.  Pulling the domain
and codomain forms through the block-diagonal basis changes therefore turns
the anti-symplectic identity for \(G\) into the following weighted one.
Explicitly, if \(N=PGQ\), with \(P=\operatorname{diag}(P_i)\) and
\(Q=\operatorname{diag}(Q_j)\), then
\[
 S^{\mathsf T}JS=(\det S)J
 \quad\Longrightarrow\quad
 P_i^{-\mathsf T}JP_i^{-1}=(\det P_i)^{-1}J,
 \qquad Q_j^{\mathsf T}JQ_j=(\det Q_j)J.
\]
Substitution in \(G^{\mathsf T}\operatorname{diag}(J)G
=-\operatorname{diag}(J)\) gives
\[
 N^{\mathsf T}\operatorname{diag}(r_0J,r_1J,r_2J)N
   =-\operatorname{diag}(s_0J,s_1J,s_2J)              \tag{5.8}
\]
with \(r_i=(\det P_i)^{-1}\) and \(s_j=\det Q_j\).  Applying the same
observation to \(N^{-1}\) gives \(t_j=s_j^{-1}\), so distinct block rows
satisfy
\[
 \sum_{j=0}^2t_jN_{ij}JN_{kj}^{\mathsf T}=0
 \qquad(i\ne k).                                     \tag{5.9}
\]
Thus \((5.8)\) gives weighted column orthogonality and \((5.9)\) gives
weighted row orthogonality.

Orthogonality of the first two block rows gives
\[
 t_0I+t_1A^{\mathsf T}+t_2B^{\mathsf T}=0,
\]
so \(B\in\operatorname{span}_{\F_q}\{I,A\}\).  The first and third rows
similarly give \(D\in\operatorname{span}\{I,C\}\).  Orthogonality of the
first two block columns gives
\[
 r_0I+r_1A+r_2C=0,
\]
so \(C\), and hence \(D\), also lies in
\(\operatorname{span}\{I,A\}\).  All four non-tree blocks commute.  In
the normalization \((5.7)\), the fundamental holonomies are their inverses,
so they commute as well.  The common centralizer is therefore at least
two-dimensional.

For \(q\geq5\), every determinant-one group attached to a nonscalar row of
Table~\ref{tab:endomorphism-types} has order strictly larger than
\(d=\gcd(2,q-1)\), which proves the final assertion.  At \(q=3\), for
example, the split norm-one group can coincide with the center; this is why
the qualification is necessary.
\end{proof}

Thus extra label-space linearity is forced through six parties, without an
occurrence or genericity claim for the remaining types.
